% Vietnamese translation for OI Public Library
% Source-File: IOI中国国家候选队论文/2008/Day1/7.方戈《浅析信息学竞赛中一类与物理有关的问题》/浅析信息学竞赛中一类与物理有关的问题.ppt
\documentclass[11pt]{article}
\usepackage{oipl-vn}
\usepackage{tikz}
\usetikzlibrary{arrows.meta,calc}

\newcommand{\Oh}{\mathcal{O}}

\title{Bản trình chiếu: Một lớp bài toán liên quan tới vật lý}
\author{Fang Ge}
\date{2008}

\begin{document}
\maketitle

\origtitle{Bản trình chiếu: Bàn về một lớp bài toán liên quan tới vật lý trong thi tin học}
\sourcepath{IOI China 2008 / Day 1 / Fang Ge / physics problems.ppt}

\section{Tinh thần của bài nói}

Slide của Fang Ge lấy vật lý làm ví dụ cho quá trình mô hình hóa. Người giải
không cần nhớ nhiều công thức vật lý nâng cao, nhưng phải biết tách hiện tượng
thành đại lượng, ràng buộc và sự kiện có thể tính được. Tài liệu gốc xoay
quanh ba bài: \textit{Ars Longa}, \textit{Water Tanks} và \textit{3D Musical
Water-fence}. Trong PPT, phần \textit{Water Tanks} chiếm dung lượng lớn nhất;
hai bài còn lại làm nền cho thông điệp chung.

\section{Từ hiện tượng sang mô hình}

Các bài kiểu này thường bắt đầu bằng một mô tả rất tự nhiên: thanh chịu lực,
nước dâng qua ống, áp suất khí thay đổi, nước chảy trên địa hình. Nếu mô phỏng
trực tiếp từng khoảnh khắc, trạng thái sẽ phình ra nhanh. Bước đầu tiên là tìm
đại lượng bất biến hoặc quan hệ cục bộ:
\begin{itemize}
  \item cân bằng lực biến thành hệ phương trình tuyến tính;
  \item khí nén tuân theo \(P_1V_1=P_2V_2\);
  \item nước không nén được, nên thể tích đã đổ bằng tổng thể tích nước đang
  ở các khoang cộng phần đã tràn;
  \item nước chảy trên địa hình có thể xử lý theo bồn trũng thay vì theo từng
  giọt nước.
\end{itemize}
Sau bước này, bài toán vật lý trở thành đại số, quét sự kiện, cấu trúc dữ liệu
hoặc nén trạng thái.

\section{\textit{Ars Longa}: cân bằng thanh}

\textit{Ars Longa} cho một hệ điểm và thanh trong không gian. Một số điểm nằm
trên mặt đất và cố định; các điểm còn lại chịu trọng lực. Thanh không khối
lượng, cứng, chỉ truyền lực dọc theo hướng của chính nó. Cần kiểm tra hệ có
cân bằng và có ổn định hay không.

Mỗi thanh được gán một ẩn lực. Với mỗi khớp tự do, tổng lực theo ba trục
\(x,y,z\) phải bằng trọng lực ngược dấu. Vì thế \(P\) khớp tự do tạo ra \(3P\)
phương trình tuyến tính. Một thanh nối hai khớp đóng góp hai vector ngược
chiều vào hai phương trình tương ứng.

\begin{center}
\begin{tikzpicture}[scale=0.95, every node/.style={font=\small}]
  \coordinate (A) at (0,0);
  \coordinate (B) at (2.0,0);
  \coordinate (C) at (0.9,1.8);
  \draw[thick] (A) -- (C) -- (B);
  \draw[-{Latex[length=2mm]},red!70,thick] (C) -- ++(0,-0.75) node[right] {\(mg\)};
  \draw[-{Latex[length=2mm]},blue!70] (C) -- ($(C)!0.45!(A)$) node[left] {\(f_1\)};
  \draw[-{Latex[length=2mm]},blue!70] (C) -- ($(C)!0.45!(B)$) node[right] {\(f_2\)};
  \foreach \p/\lab/\pos in {A/cố định/below,B/cố định/below,C/khớp/above}
    \fill (\p) circle (2pt) node[\pos] {\lab};
\end{tikzpicture}
\end{center}

Nếu hệ tuyến tính vô nghiệm, cấu trúc không cân bằng. Nếu hệ có nghiệm nhưng
ma trận hệ số không phủ được mọi vế phải có thể phát sinh, cấu trúc cân bằng
nhưng không ổn định. Về cài đặt, khử Gauss là đủ cho kích thước trong đề.

\section{\textit{Water Tanks}: bài chính của slide}

Bài \textit{Water Tanks} có \(N\) bình trụ xếp thành hàng. Bình đầu hở, các
bình sau kín. Giữa hai bình kề nhau có ống nối ở độ cao cho trước. Nước và
khí đi qua ống; nước không nén được, khí nén theo \(P_1V_1=P_2V_2\). Cần tính
lượng nước phải đổ để bình đầu đầy, với \(N\) có thể rất lớn.

\begin{center}
\begin{tikzpicture}[x=0.75cm,y=0.55cm, every node/.style={font=\small}]
  \foreach \x/\h/\w in {0/5.8/4.6,2/4.8/3.1,4/5.4/2.2} {
    \draw (\x,0) rectangle ++(1,\h);
    \fill[blue!25] (\x,0) rectangle ++(1,\w);
  }
  \draw[thick] (1,2.7) -- (2,2.7);
  \draw[thick] (3,3.6) -- (4,3.6);
  \node[below] at (0.5,-0.25) {hở};
  \node[below] at (2.5,-0.25) {kín};
  \node[below] at (4.5,-0.25) {kín};
\end{tikzpicture}
\end{center}

Mô phỏng ngây thơ sẽ thất bại vì khi nước vượt qua một ống, khoang khí có thể
bị tách ra hoặc nhập với khoang khác. Slide dẫn người nghe qua hai bước:
\begin{enumerate}
  \item xét trường hợp độ cao ống tăng dần, nơi các sự kiện xảy ra theo thứ
  tự tự nhiên từ trái sang phải;
  \item mở rộng sang trường hợp tổng quát bằng cách nhóm một đoạn bình thành
  một khối có hành vi giống trường hợp đơn điệu.
\end{enumerate}

Trong trường hợp đơn điệu, mỗi lần một khoang bị cô lập, thể tích khí còn lại
và áp suất của nó được cố định theo luật Boyle. Phần bên trái có thể xem như
đã xử lý xong, phần bên phải tiếp tục nhận nước. Mỗi bình chỉ tham gia số lần
hằng số, nên thời gian \(\Oh(N)\).

\section{Khối trong trường hợp tổng quát}

Khi độ cao ống lên xuống tùy ý, ta không thể xử lý đơn giản từ trái qua phải.
Ý tưởng của slide là gom các bình liên tiếp thành khối. Bên trong một khối,
các ống có quan hệ độ cao cho phép mô phỏng như một đơn vị; khi khối bị tách
hoặc đầy tới một ngưỡng, nó được thay bằng một trạng thái cô đọng.

Điều cần giữ trong bản dịch là lập luận chi phí: mỗi thao tác chi tiết gắn với
một bình hoặc một ống cụ thể; sau khi đã bị gom hoặc tách, phần đó không bị
xử lý lặp lại nhiều lần. Vì vậy tổng thời gian vẫn tuyến tính và bộ nhớ cũng
tuyến tính theo \(N\). Đây là mẫu rất phổ biến trong bài mô phỏng vật lý:
không theo dõi toàn bộ hệ ở độ phân giải cao, mà thu gọn các miền có cùng hành
vi.

\section{\textit{3D Musical Water-fence}}

Bài này đặt nước trên một lưới địa hình ba chiều. Mỗi thao tác đổ nước vào một
ô; nước chảy xuống thấp, tụ lại trong bồn, rồi tràn qua cửa thấp nhất. Nếu xử
lý từng thao tác theo thứ tự đề cho, một lượng nước có thể lan qua rất nhiều ô.

Quan sát đầu tiên là các thao tác đổ nước có tính giao hoán: chỉ tổng lượng
nước rơi vào từng ô mới quan trọng cho lượng tràn cuối cùng. Vì vậy ta gom
toàn bộ thao tác trước, rồi xử lý theo cấu trúc địa hình. Quan sát thứ hai là
nên nén các vùng trũng thành bồn. Mỗi bồn có cao độ tràn và dung tích:
\[
  \operatorname{cap}(B)=h(B)\cdot |B|-\sum_{u\in B}\operatorname{height}(u).
\]
Nếu lượng nước vào bồn không vượt dung tích thì không tràn; nếu vượt, chỉ phần
dư được đẩy qua cửa tràn sang bồn thấp hơn hoặc ra ngoài lưới.

\section{Bài học mô hình hóa}

Ba ví dụ cho cùng một quy trình. Trước hết phải gọi tên đúng các đại lượng vật
lý: lực, áp suất, thể tích, mực nước, dung tích. Sau đó biến chúng thành điều
kiện toán học đủ nhỏ: phương trình tuyến tính, phương trình một ẩn, quét độ
cao, hoặc cây các bồn trũng. Cuối cùng mới chọn cấu trúc dữ liệu.

Điểm dễ sai không nằm ở công thức khó, mà ở ranh giới trạng thái: khi nào khí
bị cô lập, khi nào nước tràn, khi nào một vùng có thể nén, khi nào một giả
thiết suy biến làm mô hình mất nghĩa. Vì vậy với bài liên quan vật lý, phần
kiểm chứng mô hình quan trọng không kém phần cài đặt.

\end{document}
